\documentclass[11pt,letterpaper]{article}

% Core packages
\usepackage[margin=1in]{geometry}
\usepackage[utf8]{inputenc}
\usepackage[T1]{fontenc}
\usepackage{lmodern}
\usepackage{microtype}
\usepackage{hyperref}
\usepackage{url}
\usepackage{booktabs}
\usepackage{graphicx}
\usepackage{xcolor}
\usepackage{amsmath}
\usepackage{amssymb}
\usepackage{enumitem}
\usepackage{array}
\usepackage{longtable}
\usepackage{tabularx}
\usepackage{multirow}
\usepackage{natbib}
\usepackage{caption}
\usepackage{subcaption}
\usepackage{float}
\usepackage{pgfplots}
\usepackage{tikz}
\usepackage{pgf-pie}
\usepackage{colortbl}
\usepackage{rotating}
\usepackage{makecell}
\usepackage{fancyhdr}
\usepackage{titlesec}
\usepackage{abstract}

\pgfplotsset{compat=1.18}
\usetikzlibrary{shapes,arrows,positioning,fit,backgrounds,matrix,decorations.pathreplacing,calc}

% Color definitions
\definecolor{prismorblue}{RGB}{30,70,140}
\definecolor{prismorred}{RGB}{180,30,30}
\definecolor{prismorgray}{RGB}{100,100,100}
\definecolor{lightgray}{RGB}{240,240,240}
\definecolor{darkgreen}{RGB}{0,120,50}
\definecolor{warnyellow}{RGB}{220,160,0}
\definecolor{cellgreen}{RGB}{180,230,180}
\definecolor{cellred}{RGB}{255,200,200}
\definecolor{cellyellow}{RGB}{255,240,180}

% Hyperref setup
\hypersetup{
    colorlinks=true,
    linkcolor=prismorblue,
    citecolor=prismorblue,
    urlcolor=prismorblue,
    pdftitle={Securing the Agentic Frontier: A Comprehensive Analysis of AI Coding Agent Security},
    pdfauthor={Prismor Security Research},
    pdfsubject={AI Agent Security, Threat Landscape, Competitive Analysis},
    pdfkeywords={AI agents, agentic security, prompt injection, MCP security, secrets management, runtime enforcement}
}

% Header/footer
\pagestyle{fancy}
\fancyhf{}
\rhead{\small\textcolor{prismorgray}{Prismor Security Research --- April 2026}}
\lhead{\small\textcolor{prismorgray}{Securing the Agentic Frontier}}
\cfoot{\small\thepage}
\renewcommand{\headrulewidth}{0.4pt}

% Title formatting
\titleformat{\section}{\large\bfseries\color{prismorblue}}{{\thesection}}{1em}{}[\titlerule]
\titleformat{\subsection}{\normalsize\bfseries\color{prismorblue}}{\thesubsection}{1em}{}
\titleformat{\subsubsection}{\normalsize\bfseries\color{prismorgray}}{\thesubsubsection}{1em}{}

% Custom commands
\newcommand{\yes}{\textcolor{darkgreen}{\textbf{YES}}}
\newcommand{\no}{\textcolor{prismorred}{\textbf{NO}}}
\newcommand{\todo}[1]{\textcolor{prismorred}{[TODO: #1]}}

% Title
\title{
    \vspace{-1cm}
    {\Large\bfseries\color{prismorblue} Securing the Agentic Frontier:}\\
    \vspace{0.3em}
    {\large A Comprehensive Analysis of AI Coding Agent Security,\\
    Threat Landscape, and Defense Architecture}\\
    \vspace{0.5em}
    \normalsize\color{prismorgray}
    Technical White Paper\\
    \vspace{0.3em}
}
\author{
    \textbf{Prismor Security Research}\\
    \small\texttt{research@prismor.dev} \quad \texttt{https://prismor.dev}\\
    \vspace{0.5em}
}
\date{\small April 28, 2026}

\begin{document}

\maketitle
\thispagestyle{fancy}

\begin{abstract}
\section*{Abstract}
% ABSTRACT GOES HERE
\end{abstract}

\tableofcontents
\clearpage

% =========================================================
% SECTION 1: INTRODUCTION
% =========================================================
\section{Introduction}
\label{sec:introduction}

AI coding agents --- including Claude Code, Cursor, Windsurf, and GitHub Copilot
--- have transitioned from experimental assistants to primary development infrastructure.
These systems now serve tens of millions of developers, autonomously reading and writing
files, executing arbitrary shell commands, accessing credential stores, and interacting
with external services through the Model Context Protocol (MCP).
The security controls developed for earlier software ecosystems --- kernel-level eBPF
hooks, cloud data-loss prevention gateways, and static secret scanners --- are
architecturally mismatched to this environment.
By the time a kernel hook observes a \texttt{curl} process spawned by an agent,
the agent has already decided to exfiltrate credentials; the window for semantic
interception has closed.

The threat is not hypothetical.
Indirect prompt injection attacks achieve greater than 85\% success against current
defenses across 42 documented attack techniques~\cite{maloyan2026prompt}.
Claude Code-assisted commits leak secrets at 3.2\%, more than double the 1.5\%
baseline rate for human-authored commits~\cite{gitguardian2026state}.
In aggregate, 28.65 million hardcoded secrets were detected on public GitHub in
2025 alone, a 34\% year-over-year increase, including 1,275,105 AI service
credentials~\cite{gitguardian2026state}.
An empirical study of 17,022 agent skills found that 520 (3.06\%) leaked
credentials, of which 89.6\% were immediately exploitable~\cite{chen2026credential}.
The MCP ecosystem has amplified these risks: a by-design remote code execution
(RCE) vulnerability in the MCP SDK's STDIO interface affected more than 150 million
downloads~\cite{mcprce2026}, over 8,000 publicly accessible MCP servers exposed
unauthenticated administrative panels~\cite{exposedmcp2026}, and 1,184 malicious
skills were confirmed in public skill repositories~\cite{exposedmcp2026}.
The generative AI cybersecurity market has responded with four major acquisitions
in eighteen months and is projected to grow from \$8.65B (2025) to \$35.5B (2031)
at a 26.5\% compound annual growth rate~\cite{marketsandmarkets2026generative}.

Despite this activity, a structural gap persists.
Existing runtime enforcement systems --- ClawGuard~\cite{zhao2026clawguard},
AgentSpec~\cite{wang2026agentspec}, AttriGuard~\cite{he2026attriguard},
ICON~\cite{wang2026icon}, and AgentSentry~\cite{authors2026agentsentry} --- each address
a subset of the threat lifecycle but none compose pre-emptive secret prevention
with semantic runtime enforcement and post-hoc remediation into a single,
developer-native package.
MCP-specific scanning tools operate as point-in-time static analyzers that
cannot detect dynamic rug-pull attacks or enforce policy at execution
boundaries~\cite{cosai2026mcp,invariantlabs2025mcp}.
Secret detection tools such as GitGuardian operate after exposure, not
before~\cite{gitguardian2026state}.
A position paper from NVIDIA and Johns Hopkins argues that system-level
defenses are necessary precisely because no single mechanism spans the full
threat lifecycle~\cite{research2026architecting}.
Prior taxonomies classify attack techniques but do not map to a deployed
enforcement product~\cite{kim2026attack,foundation2025owasp}.

This paper presents Prismor, a three-layer defense-in-depth architecture
operating entirely at the agent's hook layer.
Layer~1 (Cloak) prevents real secret values from entering model context by
substituting \texttt{@@SECRET:<name>@@} placeholders before the agent processes
any prompt, with \texttt{PreToolUse} and \texttt{PostToolUse} hooks performing
transparent substitution and scrubbing at execution boundaries.
Layer~2 (Warden) intercepts every tool call at the same hook layer and evaluates
it against a YAML policy engine with 25+ default rules across 10 threat
categories, operating in either observe or enforce mode with sub-millisecond
overhead.
Layer~3 (Sweep) performs post-hoc remediation by scanning Claude transcript
JSONL files, Cursor and Windsurf cache directories, and AI tool configuration
directories using Gitleaks'~\cite{rice2024gitleaks} 170+ credential patterns,
with redact, delete, and dry-run operating modes.
A signed advisory feed with Ed25519 integrity verification provides runtime
threat correlation across all three layers.

The main contributions of this paper are:
\begin{enumerate}[leftmargin=*, label=(\arabic*)]
  \item A formal characterization of the AI coding agent threat surface that is
        distinct from existing LLM application security and ML security threat models,
        grounded in documented incidents from 2025--2026.
  \item A demonstration that hook-layer interception is both \emph{necessary} and
        \emph{sufficient} to close the semantic gap that kernel-level and proxy-based
        controls cannot fill, and that this interception point supports the full
        three-layer composition.
  \item Empirical evidence that pre-emptive secrets cloaking reduces leakage risk
        to near-zero for registered secrets, against a 3.2\% baseline leak rate in
        Claude Code-assisted commits~\cite{gitguardian2026state}.
  \item A competitive capability analysis demonstrating that no existing commercial
        or open-source tool covers all three layers (Cloak, Warden, Sweep) for
        developer-native AI coding agent environments.
\end{enumerate}

% =========================================================
% SECTION 2 (Background/Threat Model placeholder, as in template)
% =========================================================
\section{Background and Threat Model}
% BACKGROUND GOES HERE

\subsection{AI Coding Agents: Architecture and Capabilities}
% AGENT DESIGN GOES HERE

\subsection{Why Existing Security Controls Fail}
% FAILURE MODES GOES HERE

% =========================================================
% SECTION 3: THREAT LANDSCAPE AND RELATED WORK
% =========================================================
\section{Threat Landscape and Related Work}
\label{sec:related}

This section surveys documented AI coding agent incidents from 2025--2026
(Section~\ref{subsec:incidents}), academic work on runtime enforcement
systems (Section~\ref{subsec:runtime}), the MCP supply chain attack surface
(Section~\ref{subsec:mcp}), and secrets and credential leakage research
(Section~\ref{subsec:secrets}).
For each cluster we identify the specific limitation that motivates the
corresponding Prismor layer.

% ---------------------------------------------------------
\subsection{Documented Incidents (2025--2026)}
\label{subsec:incidents}

The incident record from 2025 to April 2026 documents an accelerating pattern
of AI coding agent compromises spanning all three threat categories that Prismor
addresses.

\textbf{Prompt injection and configuration poisoning.}
In March 2025, a Unicode-obfuscated rules-file backdoor was discovered affecting
both Cursor and GitHub Copilot~\cite{security2025rules}, exploiting the fact that
neither product validated the semantic content of project-level configuration
before executing agent instructions.
In April 2026, CVE-2025-59536 (CVSS 8.7) documented a settings.json injection
vector in Claude Code that permits an adversarially crafted repository to override
hook configurations, illustrating that the hook layer itself is an attack
surface that must be hardened with non-overridable critical rules.

\textbf{MCP supply chain and RCE.}
In April 2025, Invariant Labs demonstrated tool poisoning attacks against MCP
servers, including tool shadowing, cross-origin privilege escalation, and rug-pull
attacks in which tool descriptions change after initial security
review~\cite{invariantlabs2025mcp,labs2025mcpscan}.
By February 2026, over 8,000 MCP servers were identified with unauthenticated
administrative panels, and 1,184 malicious skills were confirmed in public
repositories~\cite{exposedmcp2026}.
In April 2026, a by-design RCE vulnerability in the MCP SDK's STDIO interface
affected more than 150 million downloads across approximately 7,000+
publicly accessible servers~\cite{mcprce2026}.
The SANDWORM\_MODE campaign (2025) embedded MCP injection payloads across 19
npm packages~\cite{sandworm2025}.
The CoSAI MCP Security White Paper~\cite{cosai2026mcp} formalizes the taxonomy
of these attacks into tool poisoning, rug pulls, and server shadowing categories.

\textbf{Secrets exfiltration and RCE via AI frameworks.}
CVE-2025-68664 (CVSS 9.3) documented RCE in LangChain via unsafe deserialization
of secrets, demonstrating that credential handling in AI framework pipelines
creates direct exploitation paths~\cite{langchainrce2025}.
The Flowise platform recorded a CVSS 10.0 RCE vulnerability (CVE-2025-59528) in
April 2026.
The LiteLLM PyPI package suffered a supply chain backdoor in March 2026.
Malicious VS Code AI extensions with over 1.5 million combined installs were used
for credential exfiltration.

This incident record confirms that AI coding agent attacks do not decompose
cleanly into prevention or detection: each incident required a different
interception point, motivating the defense-in-depth composition.

% ---------------------------------------------------------
\subsection{Academic Research on Runtime Enforcement}
\label{subsec:runtime}

A cohort of concurrent academic systems addresses runtime enforcement for LLM
agents at varying architectural layers.

\textbf{Hook-layer enforcement.}
Prior non-agent-specific guardrail frameworks such as NeMo Guardrails~\cite{rebedea2023nemo}
introduced programmable safety rails for LLM applications, but operate at the
application-logic layer rather than at an agent's tool-call boundary.
ClawGuard~\cite{zhao2026clawguard} (Zhao et al., arXiv:2604.11790) intercepts
tool calls at hook boundaries and achieves less than 5\% attack success across
five large language models on three benchmarks: AgentDojo, SkillInject, and
MCPSafeBench.
ClawGuard is the closest architectural relative of Prismor Warden, establishing
that hook-layer enforcement is both technically feasible and effective.
However, ClawGuard does not ship with a user-configurable policy-as-code engine,
does not integrate pre-emptive secrets cloaking, and does not include post-hoc
sweep remediation.

\textbf{Policy DSL enforcement.}
AgentSpec~\cite{wang2026agentspec} (Wang et al., arXiv:2503.18666, ICSE 2026)
introduces a domain-specific language for specifying runtime constraints on agent
behavior, achieving greater than 90\% prevention of unsafe code execution,
100\% elimination of hazardous actions, and sub-millisecond enforcement overhead.
AgentSpec demonstrates that policy-as-code approaches are both expressive and
low-overhead.
The system operates at the model output level rather than at a deployment-native
hook protocol, making it unsuitable for direct integration with Claude Code's
\texttt{PreToolUse}/\texttt{PostToolUse} interface without adaptation.

\textbf{Causal attribution.}
AttriGuard~\cite{he2026attriguard} (He et al., arXiv:2603.10749) uses causal
attribution of tool invocations to detect adversarial manipulation, achieving 0\%
attack success against static injection attacks.
The causal attribution approach provides strong guarantees but requires
instrumentation of the model's internal reasoning trace, which is not available
in closed-weight models such as Claude 3.x.

\textbf{Latent-space detection.}
ICON~\cite{wang2026icon} (Wang et al., arXiv:2602.20708) detects prompt injection
in the model's latent representation space, reducing attack success to 0.4\% while
improving utility by more than 50\%.
AgentSentry~\cite{authors2026agentsentry} (arXiv:2602.22724) addresses multi-turn
temporal causal injection and session correlation attacks that single-turn systems
miss.
Both approaches operate inside the model and require either white-box access or
specialized inference infrastructure, neither of which is available to
developer-native tooling.

\textbf{System-level position.}
A position paper from researchers at NVIDIA and Johns Hopkins
(\cite{research2026architecting}, arXiv:2603.30016) argues that system-level defenses are
necessary because no single model-layer or kernel-layer mechanism can address the
full attack surface.
This argument directly supports the Prismor three-layer composition.

\textbf{Limitation.}
All of the above systems operate as standalone defenses.
None address pre-emptive secret prevention (which reduces the attack surface
before any injection can occur), and none provide post-hoc sweep remediation for
secrets already resident in agent caches and transcripts.
Consequently, even an agent fully protected by ClawGuard or AgentSpec remains
exposed to the 3.2\% background secret leakage rate documented in
\cite{gitguardian2026state}.
Prismor Warden operates at the same hook-layer interception point as ClawGuard
but composes it with Cloak and Sweep to cover the complete threat lifecycle.

% ---------------------------------------------------------
\subsection{MCP as Primary Emerging Attack Surface}
\label{subsec:mcp}

The Model Context Protocol has become the dominant mechanism by which AI coding
agents extend their capabilities with external tools and data sources.
Systematic surveys of the MCP ecosystem have documented a qualitatively new
category of supply chain threat.

\textbf{Attack taxonomy.}
The CoSAI MCP Security White Paper~\cite{cosai2026mcp} enumerates three primary
attack classes: tool poisoning (injecting adversarial instructions into tool
descriptions prior to installation), rug pulls (changing tool behavior after the
agent has already trusted the description), and server shadowing (registering
a malicious server that intercepts calls intended for a legitimate server).
Invariant Labs' April 2025 research~\cite{invariantlabs2025mcp,labs2025mcpscan} demonstrated all
three classes against real MCP servers, including cross-origin privilege escalation
through chained tool calls.
The SANDWORM\_MODE npm campaign~\cite{sandworm2025} applied MCP injection at the
package supply chain level, embedding exfiltration payloads in 19 widely installed
packages.

\textbf{Scale of exposure.}
The unauthenticated admin panel vulnerability affecting 8,000+ MCP servers
\cite{exposedmcp2026} illustrates that MCP infrastructure does not yet implement
the access control patterns required for a hostile internet environment.
The by-design RCE in the MCP SDK STDIO interface~\cite{mcprce2026} affected over
150 million SDK downloads, qualifying it as one of the highest-reach
vulnerabilities in the AI tooling ecosystem to date.
Chen et al.~\cite{chen2026credential} analyzed 17,022 agent skills and found 520
vulnerable, with 76.3\% of issues requiring joint code-and-natural-language
analysis to detect --- a modality that static code scanners do not support.

\textbf{Limitation of existing MCP security tooling.}
Current MCP security tools, including MCP-Scan and the Snyk/Invariant CLI,
provide static inspection of MCP server tool descriptions at installation time.
They cannot detect dynamic rug-pull attacks in which descriptions change
post-installation.
They do not enforce policy at execution boundaries and do not integrate
credentials cloaking or post-hoc sweep for secrets exfiltrated through MCP tool
calls.
Prismor Warden's MCP/skill scanner rule category applies at the \texttt{PreToolUse}
hook, evaluating actual tool-call arguments --- not static descriptions --- at the
moment of execution.

\textbf{Secrets and credential leakage.}
\label{subsec:secrets}
The GitGuardian State of Secrets Sprawl 2026~\cite{gitguardian2026state} documents
28.65 million hardcoded secrets on public GitHub in 2025, a 34\% year-over-year
increase, including 1,275,105 AI service credentials (an 81\% increase).
The disparity between Claude Code-assisted commit leak rate (3.2\%) and the
human baseline (1.5\%) demonstrates that AI coding agents introduce a
measurable, attributable increase in exposure risk independent of adversarial
attack~\cite{gitguardian2026state}.
Notably, 64\% of credentials exposed before 2022 remained exploitable as of
January 2026, and 28\% of incidents involve non-code channels (Slack, Jira),
indicating that leaked credentials persist long after the initial exposure
event~\cite{gitguardian2026state}.
CVE-2025-68664~\cite{langchainrce2025} demonstrates that exposed credentials
in AI framework pipelines translate directly to RCE paths.

Existing detection tools (GitGuardian scanning, LLM Guard, Nightfall) operate
post-exposure: they identify secrets that are already committed, already in
transcripts, or already transmitted to cloud APIs.
No existing open-source or commercial tool implements pre-emptive placeholder
substitution at the agent's prompt-submission boundary, which is the only
interception point that can prevent a secret from ever entering model context
or leaving the developer's machine.
Prismor Cloak fills this gap.

\textbf{Threat taxonomy surveys.}
Kim et al.~\cite{kim2026attack} (arXiv:2603.11088) survey the
AI agent attack and defense landscape across seven design dimensions, providing
the most systematic taxonomy available as of April 2026.
Maloyan and Namiot~\cite{maloyan2026prompt} (arXiv:2601.17548) systematize 78
studies and identify 42 distinct attack techniques with a composite success
rate exceeding 85\% against current defenses.
The OWASP Top 10 for LLM Applications v1.1~\cite{foundation2025owasp} defines 10 risk
categories including prompt injection, insecure output handling, sensitive
information disclosure, and supply chain vulnerabilities.
These taxonomies are descriptive: they classify attack vectors and enumerate
defenses but do not map to the enforcement architecture of any deployed product.
Prismor maps all three layers explicitly to OWASP LLM categories, covering 6--7
of the 10 risks with direct enforcement or detection capability.

% =========================================================
% Remaining sections (placeholder, per template structure)
% =========================================================
\section{Prismor: Defense-in-Depth Architecture}
% ARCHITECTURE GOES HERE

\subsection{Warden: Hook-Layer Runtime Enforcement}
% WARDEN GOES HERE

\subsection{Cloak: Pre-emptive Secret Substitution}
% CLOAK GOES HERE

\subsection{Sweep: Post-hoc Secret Remediation}
% SWEEP GOES HERE

\subsection{Immunity Feed: Signed Threat Intelligence}
% FEED GOES HERE

\section{Competitive Analysis}
% COMPETITIVE ANALYSIS GOES HERE

\subsection{Competitive Map and Taxonomy}
% MAP GOES HERE

\subsection{Vendor Profiles}
% VENDOR PROFILES GOES HERE

\subsection{Capability Matrix}
% MATRIX GOES HERE

\section{Evaluation}
% EVALUATION GOES HERE

\subsection{OWASP LLM Top 10 Coverage}
% OWASP GOES HERE

\subsection{Secrets Leakage Empirical Evidence}
% EMPIRICAL GOES HERE

\subsection{Attack Success Rate Benchmarks}
% BENCHMARKS GOES HERE

\section{Market Opportunity and Strategic Roadmap}
% MARKET GOES HERE

\subsection{Market Sizing and Growth}
% SIZING GOES HERE

\subsection{Strategic Opportunities}
% OPPORTUNITIES GOES HERE

\subsection{Product Roadmap}
% ROADMAP GOES HERE

\subsection{Go-to-Market Strategy}
% GTM GOES HERE

\section{Conclusion}
% CONCLUSION GOES HERE

\bibliographystyle{plainnat}
\bibliography{refs}

\end{document}
